This paper presented LORAD, an experimental optical proximity link, and showed
that replacing threshold decoding with maximum-likelihood sequence detection
substantially improves reliability with no change to the transmitter, the
receiver electronics, or the modulation. For a mass- and power-constrained
surface asset, that distinction is the point: the improvement is obtained
entirely from information the existing receiver already captures and discards,
and costs only computation that need not reside on the asset itself.

The central measurement is that the link's dominant failure mode is symbol
\emph{omission} rather than misclassification. Pulse-width classes remain
separated by 12--40 pooled standard deviations at every rate tested, and at the
lowest rate no substitution errors occurred across 67{,}581 symbols. What limits
the link is symbols the receiver never detects: at 114.3~bit/s, 52.5\% of
transmissions were missing at least one, and because a byte spans nine pulses, a
per-symbol loss of 0.371\% compounds into majority message failure on a
25-character message. This has a direct architectural consequence: a better
per-symbol classifier cannot help, because the missing observation was never
produced. It is not visible from aggregate error rates alone.

Every undetected symbol nonetheless leaves a signature in the inter-pulse gap,
present in 100\% of the single-omission transmissions examined. A sequence
decoder combining a learned emission model, gap-timing evidence, and the
protocol's framing and printable-ASCII constraints exploits that signature. It
raises transmission success from 72.2\% to 100\% at 28.6~bit/s and from 36.1\% to
69.4\% at 114.3~bit/s on held-out data, and in live closed-loop operation from
5.9\% to 44.9\% at 114.3~bit/s, recovering 41\% of the transmissions the
threshold decoder lost.

Framed in mission terms, roughly two in five transmissions that would have
required retransmission are instead recovered at first reception. On a
proximity link operating within scheduled communication windows and a
constrained energy budget, avoided retransmissions translate directly into
window time and energy returned to the mission, and on a store-and-forward
relay path they avoid delivery latency that can extend to an entire orbit.

Four limitations bound these results. First, the decoder recovers only 42\% to
57\% of the available headroom above the lowest rate (69.4\% against a computed
ceiling of 94.4\% at 114.3~bit/s), and its residual failures are almost
entirely single-character substitutions arising where the payload constraint
admits more than one completion. Second, sample sizes of approximately 36
held-out transmissions per rate give confidence intervals of roughly $\pm$15
percentage points; these separate the two decoders at four of the five rates but
overlap at 171.4~bit/s, and they do not resolve the shape of the per-rate
curves. Third, the held-out split contains no natural-language payloads, so the
offline figures characterize the decoder on random and adversarially framed
messages rather than on representative traffic, and the comparison between
payload families that the dataset was built to support remains unmade. Fourth,
and most consequentially for a flight application, the decoder does not fit in
the receiver's memory and runs on a host computer over the receiver's serial
link.

That third limitation is less restrictive than it first appears. The
demonstrated architecture, a constrained endpoint streaming observations to a
better-resourced decoder and receiving corrected messages back, mirrors a
plausible surface topology, in which a rover offloads decoding to a lander or
relay node with greater processing and thermal margin. Where that topology is
unavailable, a fixed-lag decoder committing decisions a bounded number of
symbols behind the observation front would bound memory and permit on-device
operation. The remaining substitution errors are exactly what a character-level
prior over the payload addresses, and the dataset deliberately mixes natural
text with random ASCII so that such a prior's contribution can be measured
rather than assumed.

Two findings are methodological rather than architectural, and are reported
because both were initially misread as properties of the channel.
Instrumentation dead time (one serial write per detected pulse) inflated
apparent symbol loss by up to a factor of two and produced
a rate-dependent degradation closely resembling channel behaviour; only
buffering the output and repeating the experiment separated the two. Separately,
because the zero and one symbols differ in duration, throughput and failure
probability both depend on the payload, and a difference in message length
between two experiments accounted for most of an apparent discrepancy in their
results. Both hazards are inherent to instrumenting a resource-constrained
receiver, and both would be harder to diagnose on a deployed system than on a
bench.

Finally, the platform's original localization objective remains open and remains
motivated: a surface asset that could derive bearing from a link it already
maintains would gain a navigation aid without additional hardware, in an
environment where GNSS is unavailable. Nothing here advances it, and one
measurement stands in the way: the transimpedance stage saturates, pinning peak
amplitude to within 0.17\% across every pulse recorded, so the present front end
yields no amplitude observable of any kind. The dual-photodiode receiver
supports a normalized-difference bearing observable that would cancel range and
gain, and the collection and analysis pipeline already accommodates angular
labels; realizing it requires reducing the transimpedance gain so the amplifier
stays linear, and angling the detectors apart, which are at present mounted
close to coaxial. Whether amplitude varies usefully with range on a linear front
end is untested: every measurement reported here was taken at a single fixed
geometry.
